\scanpage{14}
\mspage{12}

X. De plus, on a la formule\rlap{\raisebox{-.1ex}{\begin{tikzpicture}[x=.45cm,y=.45cm,baseline=0pt]\draw[line width=.4pt] (1,2.34)--(.93,2.31);\draw[line width=.4pt] (.81,2.25)--(.78,2.23);\draw[line width=.4pt] (.53,2)--(.5,1.98);\draw[line width=.4pt] (.43,1.89)--(.39,1.82);\draw[line width=.4pt] (.33,1.72) .. controls (.16,1.42) and (.27,1.23) .. (.59,1.09) .. controls (.81,.98) and (.98,1.05) .. (1.23,.99) .. controls (1.4,.94) and (1.58,.94) .. (1.71,.94) .. controls (1.87,.94) and (1.81,.72) .. (1.65,.6);\draw[line width=.4pt] (1.94,.73) .. controls (2.02,.79) and (2.09,.81) .. (2.2,.81);\draw[line width=.4pt] (2.79,.83)--(2.83,.83);\draw[line width=.4pt] (2.89,.83)--(3.33,.81);\draw[line width=.4pt] (3.4,.81)--(3.54,.82);\draw[line width=.4pt] (1.8,.32)--(1.9,.04);\end{tikzpicture}}}\inserted{\handwritten{du type $g^{\bullet}f_{\bullet}=f'_{\bullet}g'^{\bullet}$ ci-dessus, impliquant la \uline{formule}\rlap{\hspace{-0.8em}\raisebox{1.6ex}{\begin{tikzpicture}[x=.2cm,y=.2cm,baseline=0pt]\draw[line width=.4pt] (0,0) .. controls (.2,.02) and (.4,-.02) .. (.55,-.02);\draw[line width=.4pt] (.63,-.02) .. controls (.9,-.03) and (1.15,-.1) .. (1.4,-.16);\end{tikzpicture}}}}} de projection habituelle. Ceci posé, les
constructions bien connues en théorie des correspondances permettent
de définir une catégorie $\underline{B}$ dont lds objets sont ceux de C, et les
Hom étant donnés par
\[
\mathrm{Hom}_{B}(X,Y)=B(XxY)\quad .
\]
On a un isomorphisme canonique $\underline{B}\simeq\underline{B}^{o}$ , qui est l'identité sur les
objets et la symmétrie $B(XxY)\simeq B(YxX)$ sur les flêches, et un foncteur
évident \deleted{$\unsure{C}\,\unsure{:}\,\underline{B}$} $F.$ : $C\rightharpoonup\underline{B}$ , qui est l'identité sur les objets et don\unsure{n}
par les graphes sur les flêches. Composé avec la symmétrie précédente
on trouve un foncteur $F^{\bullet}$: \deleted{$\unsure{B}($} $C\rightharpoondown\underline{B}$. On vérifie enfin que le couple
($F.$,$F^{\bullet}$) satisfait aux conditions envisagées dans l'alinéa 1 (ceci aus\illegible{1}
devrait être formalisé au cas des correspondances géné rales ...). \unsure{\deleted{D}O|\deleted{L}O|\deleted{O}D|\deleted{O}L}\unsure{\deleted{a}n|\deleted{n}a}
vérifie enfin qu'il a la propriété universelle qu'il faut. Lorsqu'on
se donne un couple ($F!$,$F'^{\bullet}$) à val eursdans une catégorie B', il faut
trouver une factorisation par B. Elle est toute trouvée sur les
objets, il faut la définir sur les flêches, donc définir, pour X,Y
dans C, une application
\[
(*)\qquad B(XxY)\to\mathrm{Hom}(X',Y')\quad .
\]
Mais si Z est au\unsure{\deleted{d}-|\deleted{-}d|\deleted{a}d|\deleted{d}a}dessus de XxY, d'où $p:Z\rightharpoonup X$ et $q:Z\rightharpoondown Y$, on lui assoc
un élément du deuxième membre de (*) comme le composé
\[
X'\xrightarrow{F'^{\bullet}(p)}Z'\xrightarrow{F!(\unsure{q})}Y'\quad .
\]
Il faut prouver que l'élément obtenu en remp\unsure{\deleted{l}a|\deleted{a}l}çant Z\inserted{\begin{tabular}[b]{@{}l@{}}\handwritten{$\unsure{(=}\unsure{Z}_{\unsure{o}})$}\\[-.9ex]
\begin{tikzpicture}[x=.45cm,y=.45cm,baseline=0pt]\draw[line width=.4pt] (.28,1.83)--(.37,1.8);\draw[line width=.4pt,line cap=round] (.15,1.67)--(.16,1.67);\draw[line width=.4pt] (.12,.57)--(.06,.48);\draw[line width=.4pt,line cap=round] (.03,.38)--(.04,.38);\draw[line width=.4pt] (.02,.16)--(.05,.08);\end{tikzpicture}
\end{tabular}} par un $Z_1$ cobor\unsure{\deleted{a}d|\deleted{d}a}
dant au dessus de X\inserted{\begin{tabular}[b]{@{}l@{}}\handwritten{\unsure{e}st le \unsure{même}}\\[-.9ex]
\begin{tikzpicture}[x=.7cm,y=.7cm,baseline=0pt]\draw[line width=.4pt] (0,0) .. controls (.24,.18) and (.48,.3) .. (.7,.31) .. controls (1.15,.31) and (1.9,.22) .. (2.28,.29);\draw[line width=.4pt,line cap=round] (2.57,.34)--(2.59,.34);\draw[line width=.4pt] (2.65,.35)--(3.02,.4);\draw[line width=.4pt] (3.12,.43)--(3.85,.56)--(3.77,.67);\draw[line width=.4pt,line cap=round] (4.08,.55)--(4.1,.55);\end{tikzpicture}
\end{tabular}}. Or si $g$: $T\rightharpoonup XxY$ réalise le cobordisme, on
voit (regardant seulement T et la décomposition bord(T) = $Z_{\handwritten{\unsure{o|0}}}\amalg Z_1$)
qu'il existe une variété $\overline{T}$ sur $S^1$ (obtenue en ``doublant'' T) \deleted{et une}
\deleted{application T \unsure{une}} tr\unsure{\deleted{q}a|\deleted{a}q}nsversale au dessus d'un couple de poihts $s_o$,s
de $S^1$, tels que $Z_i$ soit l'image inverse \inserted{\begin{tabular}[b]{@{}l@{}}\handwritten{dans $\overline{T}$}\\[-.9ex]
\begin{tikzpicture}[x=.68cm,y=.68cm,baseline=0pt]\draw[line width=.4pt] (3.25,1.65) .. controls (3.34,1.4) and (3.26,1.27) .. (3.12,1.25) .. controls (2.5,1.18) and (1.45,1.14) .. (.73,1.06) .. controls (.4,1.05) and (.12,.91) .. (.05,.63) .. controls (-.02,.36) and (.07,.17) .. (.25,0);\end{tikzpicture}
\end{tabular}} de $s_i$ (i=0,1); d'autre part,

\margin{13}

